\renewcommand{\baselinestretch}{1.5}

\chapter{Hypothesis and Objectives}

\section{Problem}

La generación de imágenes sintéticas con modelos de difusión latente se ha estudiado sobre todo con conjuntos de datos grandes y diversos. En agricultura de precisión, en cambio, los datos suelen ser pocos, homogéneos (las imágenes se parecen mucho entre sí: mismo cultivo, altura de vuelo y tipo de escena, con poca variabilidad) y, cada vez más, multimodales (cada imagen combina varias fuentes de información, por ejemplo color y profundidad). Este trabajo parte de un conjunto de 4\,908 imágenes aéreas de plantaciones de piña (resolución 960$\times$544), con una variante que añade un canal de profundidad monocular estimada.

Un modelo de difusión latente se entrena en dos etapas: un autoencoder variacional (VAE) comprime la imagen a un espacio latente (aquí, 4$\times$32$\times$32 para recortes de 256$\times$256) y un generador (DDPM o Flow Matching) aprende a producir latentes a partir de ruido. La calidad del espacio latente fija el techo de lo que el generador puede lograr, pero las decisiones sobre qué datos alimentan al VAE y cómo se estructura se toman hoy de forma manual y su efecto sobre la \emph{generación} no está claro. En particular, la literatura reciente señala que mejorar la reconstrucción de un VAE no implica mejorar la generación, pero esto se ha estudiado con grandes volúmenes de datos y no en un régimen de pocos datos homogéneos.

El problema que aborda esta propuesta es determinar cómo la variación de los datos de entrenamiento y de la estructura del VAE afecta la calidad de las imágenes aéreas de piña generadas mediante difusión latente, y a qué costo.

\section{Hypothesis}

Se plantean dos hipótesis, una sobre el efecto de manipular los datos y la estructura del VAE y otra sobre el límite de esa mejora:

\begin{enumerate}
  \item \textbf{Mejora.} Aplicar tratamientos a los datos de entrenamiento del VAE (aumento, enmascarado aleatorio, selección de imágenes y modalidad, con o sin canal de profundidad) y modificar su estructura (capa variacional y pérdida adversarial con un discriminador tipo GAN) mejorará el FID y el KID de las imágenes aéreas de piña generadas con modelos de difusión (DDPM y Flow Matching), en comparación con un VAE base sin tratamientos (entrenado solo con recortes aleatorios de 256$\times$256), sobre un conjunto de datos pequeño, homogéneo y multimodal.
  \item \textbf{Límite.} La mejora de la generación al acumular estos tratamientos y modificaciones presentará rendimientos decrecientes: existirá un punto a partir del cual aumentar la calidad de reconstrucción del VAE no reducirá significativamente el FID ni el KID.
\end{enumerate}

La primera hipótesis se refuta si ninguna combinación de tratamientos y modificaciones reduce el FID/KID de forma significativa respecto al VAE base; la segunda, si el FID/KID sigue reduciéndose significativamente en todos los escalones de la manipulación.

\section{Objectives}

\subsection{Main Objective}

Evaluar experimentalmente el efecto de variar los datos de entrenamiento y la estructura del VAE sobre la calidad de las imágenes aéreas de piña generadas mediante difusión latente, en un régimen de pocos datos homogéneos y multimodales.

\subsection{Specific Objectives}

\begin{enumerate}
  \item Implementar una línea base reproducible del pipeline VAE + generador (DDPM y Flow Matching), con una receta de entrenamiento controlada y verificada.
  \item Comparar tratamientos de los datos de entrada del VAE: aumento, enmascarado aleatorio, selección de imágenes y modalidad (RGB y RGB con profundidad).
  \item Comparar modificaciones estructurales del VAE: capa variacional, pérdida adversarial (GAN) y regularización.
  \item Analizar estadísticamente, con varias semillas, la relación entre la calidad de reconstrucción del VAE y la calidad de generación (FID/KID).
  \item Experimentar con un encoder auto-supervisado preentrenado (DINO) como espacio latente alternativo.
  \item Elaborar un artículo con los resultados del estudio.
\end{enumerate}

\subsection{Deliverables}

\begin{table}[h]
\centering
\small
\begin{tabular}{p{0.34\textwidth}p{0.58\textwidth}}
\hline
\textbf{Objetivo específico} & \textbf{Entregable} \\
\hline
1. Línea base reproducible & Informe de auditoría del pipeline y de la receta de entrenamiento corregida; código y configuraciones. \\
2. Efecto de los datos & Familia de VAEs entrenados con cada tratamiento de los datos; tabla comparativa de FID/KID y reconstrucción. \\
3. Efecto de la estructura & VAEs con capa variacional, GAN y regularización; análisis de su efecto sobre la reconstrucción y la generación. \\
4. Reconstrucción contra generación & Análisis estadístico (correlaciones, curvas, intervalos de confianza) con varias semillas. \\
5. Espacios latentes alternativos & Prototipo con un encoder auto-supervisado (DINO) y su comparación con el VAE base. \\
6. Artículo & Manuscrito con los resultados del estudio. \\
\hline
\end{tabular}
\caption{Entregables asociados a cada objetivo específico.}
\end{table}

\section{Research questions}

\begin{itemize}
  \item \textbf{RQ1.} ¿Cómo afecta la variación de los datos de entrenamiento del VAE (tratamientos: aumento, enmascarado, selección, modalidad) a la calidad de las imágenes generadas?
  \item \textbf{RQ2.} ¿Cómo afectan las modificaciones estructurales del VAE (capa variacional, GAN, regularización) a la calidad de las imágenes generadas?
  \item \textbf{RQ3.} ¿Se traduce una mejor reconstrucción del VAE en una mejor generación, o existe un punto a partir del cual deja de ayudar?
  \item \textbf{RQ4.} ¿Puede un encoder auto-supervisado preentrenado sustituir a un VAE entrenado desde cero en este régimen de datos?
\end{itemize}

\section{Scope and Limitations}

\subsection{Scope}

\begin{itemize}
  \item Un único conjunto de datos de imágenes aéreas de piña (4\,908 imágenes; 3\,536 para entrenamiento, 381 para validación y 346 para prueba), con recortes de 256$\times$256.
  \item VAE convolucional de espacio latente 4$\times$32$\times$32 y dos generadores: DDPM y Flow Matching, ambos con una U-Net convolucional sin atención.
  \item Calidad medida con FID y KID sobre imágenes generadas y un conjunto de referencia fijo; reconstrucción medida con MSE y PSNR.
  \item Núcleo del estudio: objetivos 1 a 4. Extensión, si el tiempo lo permite: objetivo 5.
  \item Extensión: la validación en otros dominios con conjuntos de datos pequeños y homogéneos, como las imágenes médicas, queda como extensión del estudio y no forma parte del núcleo.
\end{itemize}

\subsection{Limitations}

\begin{itemize}
  \item No se evalúa la utilidad de las imágenes sintéticas en una tarea posterior (por ejemplo, detección); se mide solo su calidad generativa.
  \item La profundidad es una estimación monocular relativa, no una medición de sensor.
  \item El número de semillas por condición está acotado por el costo de cómputo (unas 28 horas de GPU por entrenamiento).
  \item El FID se calcula con un número limitado de imágenes generadas (400), por lo que se acompaña de KID y de intervalos de confianza.
  \item Los resultados corresponden a un cultivo y a un conjunto de datos; su generalización a otros dominios no se evalúa.
\end{itemize}
